\section{Conclusion}
\label{sec:conclusion}

We surveyed over 90 memory systems for LLM-based agents and organized
them along four axes: structure, retrieval, lifecycle, and preference
emergence. Three findings stand out.

First, the AI agent community has independently converged on five
design principles---ranked retrieval, typed stores, graph traversal,
temporal decay, and automatic extraction---that were predicted by both
cognitive neuroscience and Yogac\=ara Buddhist psychology centuries or
decades before LLMs existed. This convergence suggests these are
not arbitrary design choices but structural necessities of any
sufficiently complex memory system.

Second, principled forgetting, preference emergence, and cross-memory
consolidation remain critically underserved. Most systems treat memory
as append-only storage with vector retrieval---the equivalent of
never forgetting and never learning from experience. The few systems
that implement lifecycle processes (SYNAPSE, LightMem, Alaya) achieve
measurably better retrieval quality.

Third, the Yogac\=ara framework provides a useful vocabulary for
discussing memory architecture that complements neuroscience models.
Where neuroscience offers mechanistic models (Hebbian LTP, CLS replay),
Yogac\=ara offers phenomenological insight: memory is perspective-relative
(\emph{vijnapti-m\=atra}), preferences emerge from accumulated traces
(\emph{v\=asan\=a}), and the storehouse itself must periodically
transform (\emph{\=a\'sraya-par\=av\d{r}tti}).

Memory is not a database problem. It is a process problem. The systems
that recognize this---and design accordingly---define the next
generation of AI agent memory.
